\chapter{Linear Algebra over $\mathbb{Q}$}\label{ch:linalg}

Row reduction was in the engine from the first week, as a step-recording algorithm driven by the
simplifier, and for months its steps were labelled \status{unverified}. This chapter is about the
two theorems that changed the label: elimination preserves the solution set, and the result is in
reduced row echelon form. Both live in the engine itself, with no Mathlib, and the second cost a
lemma that reads better than the algorithm it is about.

\section{A matrix is a system}

The first decision was what a matrix \emph{means}. Reading it as a linear map would have needed
matrix multiplication and its algebra; reading it as a homogeneous system needs a dot product.
One equation $r \cdot x = 0$ per row, and the augmented matrix of $Ax = b$ is the same system at
the point $(x, -1)$, so both readings are one predicate:

\begin{minted}{lean}
def Sol (m : Mat) (x : List Rat) : Prop := ∀ r ∈ m, dot r x = 0
\end{minted}

Elimination is a list of the three elementary row operations, and the theorem is one lemma per
operation folded over the list:

\begin{minted}{lean}
inductive Op where
  | swap (i j : Nat)              -- exchange rows i and j
  | scale (i : Nat) (c : Rat)     -- Rᵢ ← c · Rᵢ; with c = 0 it is the identity
  | addMul (i j : Nat) (c : Rat)  -- Rᵢ ← Rᵢ + c · Rⱼ; with i = j it is the identity

theorem sol_rref (m : Mat) (x : List Rat) : Sol (rref m) x ↔ Sol m x
\end{minted}

\section{Three ways to have no side condition}

The lemmas have no hypotheses, and each absence was arranged.

\begin{enumerate}[nosep]
  \item \textbf{Degenerate parameters are the identity.} Scaling a row by $0$ would lose an
    equation; adding a row to itself with coefficient $-1$ would too. Instead of proving that
    the algorithm avoids these cases, the operations \emph{define} them as doing nothing. Then
    every operation is invertible for free, and the algorithm needs no proof that it stays
    away from the edge.
  \item \textbf{Pad, do not truncate.} Adding a shorter row to a longer one pads the shorter with
    zeros (\lc{rowAdd}). This removes the rectangularity hypothesis from every lemma, because
    nothing ever depends on rows having equal length.
  \item \textbf{The field algebra is \lc{grind}'s.} Mathlib's \lc{ring} is not available in the
    Init-only engine, and the identities over \lc{Rat} that each lemma needs are discharged by
    \lc{grind}, whose ring and field instances ship with core. So the whole proof lives in
    \file{engine/}, next to the code it is about.
\end{enumerate}

The command's numeral path replays the emitted operations into the derivation steps, so the
notebook shows \rn{la.row-swap}, \rn{la.row-scale} and \rn{la.row-add} with the theorem names.
Symbolic entries keep the old simplifier-driven algorithm under a \texttt{.symbolic} suffix and
stay \status{unverified}, because a symbolic pivot is only \emph{assumed} nonzero.

\section{Echelon form: proved, not checked}

The first version decided at run time that the output was in echelon form and refused otherwise
--- the boundary of Chapter~\ref{ch:order}, applied once more. It was on the open list because a
boundary is a promise to come back. The proof carries an invariant column by column:

\begin{minted}{lean}
/-- Rows `< p` are pivot rows with pivot columns `pivs`; rows `≥ p` are zero left of `col`. -/
structure Inv (m : Mat) (p col : Nat) (pivs : List Nat) : Prop where
  len    : pivs.length = p
  lt     : ∀ c ∈ pivs, c < col
  sorted : pivs.Pairwise (· < ·)
  pivot  : ∀ i, i < p → entry m i (pivs.getD i 0) = 1
  lead   : ∀ i, i < p → ∀ j, j < pivs.getD i 0 → entry m i j = 0
  alone  : ∀ i, i < p → ∀ k, k ≠ i → entry m k (pivs.getD i 0) = 0
  below  : ∀ i, p ≤ i → ∀ j, j < col → entry m i j = 0
  p_le   : p ≤ m.length

def IsRref (m : Mat) (w : Nat) : Prop := ∃ p pivs, Inv m p w pivs
theorem rref_isRref (m : Mat) : IsRref (rref m) (ncols m)
\end{minted}

A column step either finds no pivot --- the invariant moves one column right --- or swaps, scales
and clears, giving one more pivot row. The lemma that made the proof entry-wise rather than
operational is \lc{entry_clears}: the clearing step reads all its factors from the matrix
\emph{before} any clearing, so its sequential application equals the simultaneous one, and each
row changes by its own multiple of the pivot row, which is itself untouched. Without that
observation the proof would have had to follow the loop; with it, every clause of \lc{Inv} is
checked on one entry at a time.

The statement is in the width of the \emph{input}. That sidesteps proving that row operations
preserve rectangularity; for the rectangular matrices the parser admits, the two widths coincide.

\begin{logbox}[Still open]
  Determinants and matrix products over $\mathbb{Q}$, and a value for matrices in the real
  semantics --- which is what would give \rn{diff.matrix} content. None of it is hard; all of it
  is listed, with the honest label, in Appendix~\ref{app:rules}.
\end{logbox}

\begin{trybox}
\begin{minted}{text}
rref([1,2,3;4,5,6;7,8,10])
rref([a,1;1,a])
\end{minted}
  The first is verified all the way down. The second is symbolic: its steps carry the
  \texttt{.symbolic} suffix and say which pivot was assumed nonzero.
\end{trybox}

\section{Exercises}

\begin{exercisebox}[Identity by definition]
  Show that \lc{sol_scale} would need the hypothesis $c \neq 0$ if scaling by $0$ zeroed the
  row, and that the algorithm as written never scales by $0$ anyway. Which of the two is the
  cheaper thing to know?
\end{exercisebox}

\begin{exercisebox}[The width]
  Give a non-rectangular \lc{Mat} for which \lc{IsRref (rref m) (ncols m)} holds but the
  intuitive ``reduced echelon form'' does not, and explain why the parser never produces it.
\end{exercisebox}
